\chapter{Dokumentenauswahl}
\label{ch:dokumentenauswahl}

Da Speicherkapazität in der Regel nur begrenzt verfügbar ist und eine Kopie der 
gesamten englischsprachigen Wikipedia einen Speicherbedarf von über 100~GB 
darstellt~\cite{}, wird nur ein Teil der Artikel ausgewählt, der dem zuvor 
konfigurierten Speicherbudget genügt. Dieses ist für den produktiven Indexierungslauf 
in \texttt{config/server.yaml} unter \texttt{selection.storage\_budget\_mb} festgelegt 
(aktuell 13.000~MB), sodass der 
Scoring- beziehungsweise Selection-Algorithmus zum Einsatz kommt. Dieser löst das 
Problem in zwei Teilen: zunächst wird für jeden Artikel ein Score berechnet, darauf 
aufbauend erfolgt eine budgetbeschränkte Auswahl.

\section{Metadatenextraktion}

Um aus den Wikipedia-Artikeln, welche in Form von HTML-Dokumenten innerhalb
einer ZIM-Datei vorliegen, Metadaten extrahieren zu können, müssen diese
Dokumente zunächst vorverarbeitet beziehungsweise normalisiert und bereinigt
werden.

Hierzu existiert innerhalb der Pipeline das Modul \texttt{Analysis}, welches
genau diese Aufgabe übernimmt. Resultierend werden artikelweise folgende
Metadaten extrahiert:

\begin{itemize}
    \item Titel
    \item Byte-Größe des Dokuments
    \item Anzahl der Wörter
    \item Anzahl der ausgehenden Links
    \item Liste der ausgehenden Links
    \item Anzahl der Kapitel
    \item Titel der Kapitel
    \item Summe der Aufrufe über 24 Stunden
\end{itemize}

Die Aufrufszahlen der einzelnen Artikel erhalten wir durch die Aggregation von, durch Wikimedia (QUELLE) bereitgestellten, Pageview-Dumps, welche
über einen Zeitraum eines vorher konfigurierten Tages gesammelt und Anhand der Artikeltitel, welche im vorhinein normalisiert wurden, eindeutig
zugeordnet werden können.
Zusätzlich zu den oben aufgeführten Metadaten aggregieren wir die Anzahl der Aufrufe eines Artikels an einem vorher konfigurierten Tag mithilfe von Wikimedia-Pageview-Dumps, 
welche mithilfe von vorher normalisierten Titeln mit den eigentlichen Artikeln verknüpft werden können.

ÜBERLEITUNG:
Die einzelnen Metadaten können anhand einer Konfigurationsdatei verschieden Gewichtet werden, um je nach Priorisierung unterschiedliche Ergebnisse im Scoring zu erzielen.

\section{Link-Frequenz-Map und Zwei-Pass-Architektur}


\section{Scoring-Modell}
Um Artikel iinnerhalb des Speicherbudgets priorisieren zu können, wird jedem Artikel ein einzelner,
kombinierter Score zugewiesen. Dieser Score setzt sich aus fünf Metriken zusammen, die jeweils einen anderen
Aspekt der Relevanz eines Artikels abbilden: der Umfang des Inhalts (\texttt{word\_count}), die ausgehende
Vernetzung (\texttt{link\_count}), eine Näherung an die korpusweite Wichtigkeit über die Häufigkeit,
mit der die verlinkten Zielartikel selbst referenziert werden (\texttt{importance}), die eingehende Popularität
innerhalb des Korpus (\texttt{incoming\_links}) sowie externe Aufrufzahlen (\texttt{pageviews}).

Da diese Signale unterschiedliche Wertebereiche und Verteilungen aufweisen, insbesondere \texttt{importance} als 
Summe über alle ausgehenden Links kann bei stark vernetzten Artikeln deutlich größere Werte annehmen als die übrigen Signale,
wird jedes Signal vor der Gewichtung logarithmisch gedämpft ($\log(1+x)$). Dadurch wird verhindert, dasss einzelne Ausreißerwerte,
etwa ein extrem häufig verlinkter Hub-Artikel, den Gesamtscore dominieren. Der Score eines Artikels $a$ berechnet sich
damit als gewichtete Summe:

\[
\mathrm{Score}(a)
= \sum_{s \in S} w_s \cdot
\log\left(1 + \mathrm{value}_s(a)\right)
\]

Wobei $S$ die fünf Signale und $w_s$ die zugehörigen, konfigurierbaren Gewichte bezeichnet.

Für die Berchnung von \texttt{importance} und \texttt{incoming\_links} wird vorab eine korpusweite
Link-Häufigkeitsmap benötigt, die zählt, wie oft jeder Artikeltitel als Ziel eines Links vorkommt. Diese
kann erst vollständig vorliegen, nachdem der gesamte Korpus einmal durchlaufen wurde. Das Scoring erfordert
somit zwangsläufig zwei Durcläufe, bevor die eigentliche Auswahl beginnen kann (siehe \todo{Verweis auf K6 Streaming-Architektur}).

Die Gewichte $w_s$ siind nicht fest codiert, sondern über \texttt{config/prod.yaml} konfigurierbar (Tabelle~\ref{tab:scoring-weights}). 
Somit ist die Priorisierung eine dokumentierte Design-Entscheidung.
\todo{Tabelle mit echten gewichten und messwerten}



\section{Budget-constrained Selection}
Das in der Aufgabenstellung definierte Ziel, eine Dokumentenauswahl anhand eines gegebenen Speicherbudget
umzusetzen, haben wir als 0/1-Knapsack-Problem modelliert: Jeder Artikel besitzt einen, durch die 
Scoring-Funktion berechneten 'Wert' und ein 'Gewicht' (Byte-Größe \texttt{html\_size\_bytes}). Ziel ist die 
Auswahl einer Teilmenge, sodass der Gesamtwert maximiert wird, ohne aber das Gewichtslimit (hier Speicherbudget)
zu überschreiten.

Das 0/1-Knapsack-Problem ist in einer exakten Form NP-hart, was für unsbedeutet, dass eine optimale Lösung bei einem
Korpus dieser Größenordnung nicht praktikabel berechenbar ist. Allerdings ist eine Approximation durch
Nutzung der Greedy-Heuristik möglich. Hierbei wird für jeden Artikel das Verhältnis von Score zu 
Byte-Größe ('Score pro Byte') berechnet, die Artikel werden absteigend danach sortiert, und in dieser Reihenfolge
ausgewählt, bis das Budget ausgeschöpft ist. Diese Heuristik liefert zwar keine nachweislich optimale
Lösung, ist aber als Standardansatz für Knapsack-Instanzen dieser Größe etabliert und in $O(n \log n)$
effizient berechenbar.

Das Ergebnis der Auswahl umfasst neben der eigentlichen Artikelmenge auch Kennzahlen zur Budgetausschöpfung
(\texttt{used\_mb}, \texttt{utilization}) — \todo{konkrete Werte aus dem Server-Lauf: Anzahl ausgewählter Artikel, tatsächlich genutztes Budget in }%}.
